
\documentclass[11pt,a4paper,oneside]{book}

\usepackage[T1]{fontenc}
\usepackage[utf8]{inputenc}
\usepackage{libertinus}
\usepackage{microtype}
\usepackage[a4paper,margin=27mm,headheight=15pt]{geometry}
\usepackage{amsmath,amssymb,amsthm,mathtools,bm,mathrsfs}
\usepackage{booktabs,longtable,array,tabularx}
\usepackage{enumitem}
\usepackage{xcolor}
\usepackage[most]{tcolorbox}
\usepackage{listings}
\usepackage{fancyhdr}
\usepackage{hyperref}
\usepackage[nameinlink,noabbrev]{cleveref}

\definecolor{deepblue}{HTML}{234A73}
\definecolor{softblue}{HTML}{EDF4FA}
\definecolor{deepolive}{HTML}{536B3D}
\definecolor{softolive}{HTML}{F2F6ED}
\definecolor{deepred}{HTML}{8A3A3A}
\definecolor{softred}{HTML}{FCF1F1}
\definecolor{warmgray}{HTML}{6B625C}
\definecolor{softgray}{HTML}{F5F4F2}

\hypersetup{
  colorlinks=true,
  linkcolor=deepblue,
  citecolor=deepolive,
  urlcolor=deepred,
  pdftitle={Double-Factorial Transseries and Their Inversion},
  pdfauthor={A self-contained coefficient-level research synthesis}
}

\pagestyle{fancy}
\fancyhf{}
\fancyhead[L]{\small Double-factorial transseries}
\fancyhead[R]{\small\ifnum\value{chapter}>0 Chapter~\thechapter\fi}
\fancyfoot[C]{\thepage}
\setlength{\parindent}{1.2em}
\setlength{\parskip}{0.15em}
\allowdisplaybreaks

\newtheorem{theorem}{Theorem}[chapter]
\newtheorem{proposition}[theorem]{Proposition}
\newtheorem{lemma}[theorem]{Lemma}
\newtheorem{corollary}[theorem]{Corollary}
\newtheorem{definition}[theorem]{Definition}
\newtheorem{remark}[theorem]{Remark}
\newtheorem{example}[theorem]{Example}
\newtheorem{algorithm}[theorem]{Algorithm}
\newtheorem{convention}[theorem]{Convention}

\newtcolorbox{resultbox}[1][]{
  enhanced,breakable,colback=softblue,colframe=deepblue,
  boxrule=0.7pt,arc=1mm,left=2mm,right=2mm,top=1.5mm,bottom=1.5mm,
  title={Main result},fonttitle=\bfseries,#1
}
\newtcolorbox{methodbox}[1][]{
  enhanced,breakable,colback=softolive,colframe=deepolive,
  boxrule=0.7pt,arc=1mm,left=2mm,right=2mm,top=1.5mm,bottom=1.5mm,
  title={Coefficient mechanism},fonttitle=\bfseries,#1
}
\newtcolorbox{warningbox}[1][]{
  enhanced,breakable,colback=softred,colframe=deepred,
  boxrule=0.7pt,arc=1mm,left=2mm,right=2mm,top=1.5mm,bottom=1.5mm,
  title={Interpretive warning},fonttitle=\bfseries,#1
}
\newtcolorbox{summarybox}[1][]{
  enhanced,breakable,colback=softgray,colframe=warmgray,
  boxrule=0.6pt,arc=1mm,left=2mm,right=2mm,top=1.5mm,bottom=1.5mm,
  title={Summary},fonttitle=\bfseries,#1
}

\lstset{
  basicstyle=\ttfamily\small,
  columns=fullflexible,
  frame=single,
  breaklines=true,
  showstringspaces=false,
  keywordstyle=\bfseries,
  commentstyle=\itshape,
  xleftmargin=1em,
  xrightmargin=1em
}

\newcommand{\R}{\mathbb R}
\newcommand{\C}{\mathbb C}
\newcommand{\Q}{\mathbb Q}
\newcommand{\N}{\mathbb N}
\newcommand{\e}{\mathrm e}
\newcommand{\ii}{\mathrm i}
\newcommand{\W}{\operatorname W}
\newcommand{\Log}{\operatorname{Log}}
\newcommand{\Res}{\operatorname*{Res}}
\newcommand{\coeff}[2]{\left[#1\right]#2}
\newcommand{\Ord}{\mathrm O}
\newcommand{\littleo}{\mathrm o}
\newcommand{\eps}{\varepsilon}
\newcommand{\Bell}{\mathcal B}
\newcommand{\etaD}{\eta}
\newcommand{\DF}{\mathcal D}

\title{\Huge\bfseries Double-Factorial Transseries and Their Inversion\\[3mm]
\Large Bernoulli--Bell Coefficients, Lambert-\(W\) Cores,\\
Parity Sectors, Borel Structure, and Discrete Inverses}
\author{A self-contained coefficient-level research synthesis}
\date{September 2026}

\begin{document}

\frontmatter
\hypersetup{pageanchor=false}
\maketitle
\hypersetup{pageanchor=true}

\chapter*{Abstract}
\addcontentsline{toc}{chapter}{Abstract}

The double factorial is deceptively close to the ordinary factorial, yet its
asymptotic inversion contains three structural features that must be separated:
a parity sector, a power--logarithmic growth law, and a divergent Bernoulli
correction series.  This article gives an all-orders treatment of all three.

For the parity label \(\eps\in\{0,1\}\), define
\[
 D_\eps(x)=C_\eps 2^{x/2}\Gamma\!\left(\frac{x}{2}+1\right),
 \qquad
 C_\eps=\left(\frac{2}{\pi}\right)^{\eps/2}.
\]
Then \(D_0(n)=n!!\) for even \(n\), while \(D_1(n)=n!!\) for odd \(n\).
The centered coordinate \(u=x+1\) gives the exact normal form
\[
 \log D_\eps(x)
 =\frac{u}{2}(\log u-1)+A_\eps+\Omega(u),
 \qquad
 A_0=\frac12\log\pi,\quad A_1=\frac12\log2,
\]
where the same correction \(\Omega\) serves both parity sectors.  We prove the
Borel--Laplace representation
\[
 \Omega(u)=\int_0^\infty \e^{-us}
 \frac{s/\sinh s-1}{2s^2}\,ds
\]
and obtain the complete logarithmic expansion
\[
 \Omega(u)\sim\sum_{k\ge1}\frac{c_k}{u^{2k-1}},
 \qquad
 c_k=
 \frac{(1-2^{2k-1})B_{2k}}{2k(2k-1)}
 =(-1)^k\frac{(2k-2)!\eta(2k)}{\pi^{2k}}.
\]
An exact signed remainder formula and a simple explicit remainder bound follow
from a Mittag--Leffler expansion of the Borel kernel.  Exponentiation is handled
to arbitrary order by complete Bell polynomials, yielding closed partition
sums and a triangular recurrence for every multiplicative coefficient.

For inversion, put
\[
 R_\eps=\log y-A_\eps,\qquad
 v=\W_0\!\left(\frac{2R_\eps}{\e}\right),\qquad
 X=\frac{2R_\eps}{v}=\e^{v+1},\qquad
 \Lambda=v+1=\log X.
\]
The exact Lambert core solves
\(\tfrac12X(\log X-1)=R_\eps\).  The inverse has the universal expansion
\[
 D_\eps^{-1}(y)
 \sim X-1+\sum_{n\ge1}
 \frac{d_{2n-1}(\Lambda)}{X^{2n-1}},
\]
where parity occurs only through \(R_\eps\).  We derive: (i) a triangular
coefficient-extraction recurrence; (ii) a scalar formula in ordinary partial
Bell polynomials; and (iii) a nonrecursive Lagrange--B\"urmann formula involving
the compositional inverse of the Bernoulli series.  In particular,
\(d_{2n-1}(\Lambda)\in\Lambda^{-1}\Q[\Lambda^{-1}]\), its largest possible
inverse power is \(\Lambda^{-(2n-1)}\), and its leading term is
\(-2c_n/\Lambda\).  The Lambert core is then expanded into iterated logarithms
with unsigned Stirling cycle numbers, and a second Bell-polynomial extraction
formula gives every coefficient of the resulting polynomial--logarithmic
transseries.

The article also treats the exact generalized inverse of the discrete sequence,
the standard all-integer analytic interpolation and its oscillatory inverse,
complex inverse sheets, optimal truncation and Borel singularities, symbolic
algorithms, numerical checks, and the extension of the entire calculus to
multifactorials.

\tableofcontents

\chapter*{Reader's guide}
\addcontentsline{toc}{chapter}{Reader's guide}

The main forward theorem is \cref{thm:forward-log}; the multiplicative
coefficient formulas are in \cref{thm:forward-bell}.  The inverse normal form is
derived in \cref{thm:inverse-master}; its coefficient recurrence and closed
Lagrange formula are \cref{thm:inverse-recursion,thm:lagrange-closed}.
The complete iterated-logarithm expansion is encoded in
\cref{thm:outer-complete}.  Readers concerned with the inverse of the
\emph{integer sequence}, rather than an analytic parity branch, should also
read \cref{ch:discrete}.  The distinction between parity-resolved inverses and
the inverse of a single oscillatory interpolation is explained in
\cref{ch:interpolation}.

\mainmatter

\chapter{Introduction and statement of the problem}

\section{Why the word ``inverse'' needs a definition}

For \(n\in\N\), the double factorial is
\[
 n!!=n(n-2)(n-4)\cdots,
 \qquad 0!!=(-1)!!=1.
\]
The sequence begins
\[
 1,1,2,3,8,15,48,105,384,945,\ldots
\]
and is OEIS A006882 \cite{OEIS-A006882}.  There are three different objects that
are routinely called an inverse double factorial.

\begin{enumerate}[label=\textup{(\roman*)}]
\item One may invert a smooth continuation of the even subsequence.
\item One may invert a smooth continuation of the odd subsequence.
\item One may ask for the least integer \(n\) such that \(n!!\ge y\).
\end{enumerate}
A fourth possibility is to choose one analytic function that interpolates both
parities and invert that function.  Such an interpolation is not unique, and
its inverse generally differs from both parity-resolved inverses by a correction
larger than the first Bernoulli correction.  Any complete answer must therefore
state which inverse is being expanded.

This article first solves the parity-resolved problem, which is canonical for
the asymptotics of the sequence.  It then converts those answers into an exact
formula for the discrete generalized inverse.  Finally, it analyzes the
all-integer interpolation recorded in OEIS and exhibits the extra oscillatory
transseries generated by its periodic multiplier.

\section{Relation to coefficient calculus and earlier inversion work}

The derivation is organized around the coefficient operations emphasized in
the ProveIt \emph{Combinatorial Coefficient Calculus}: exponentiation by complete
Bell polynomials, composition by partial Bell polynomials, and reversion by
Lagrange--B\"urmann extraction \cite{ProveIt-CCC}.  The Lambert-core strategy is
also aligned with the special-function inversion program in the repository's
series-and-transseries collection: invert the dominant power--logarithmic phase
exactly, then generate correction blocks above that core
\cite{ProveIt-Series,ProveIt-SpecialInverse}.  Here these ideas are specialized,
proved, and sharpened for the double factorial.  In particular, the centered
double-factorial correction has an unusually simple Borel transform, allowing
more precise remainder and late-term information than a purely formal
calculation provides.

\section{What ``full transseries'' means here}

We distinguish three levels.

\begin{definition}[Algebraic all-orders transseries]
An expansion is algebraically complete if every inverse-power and iterated-log
coefficient is specified by a finite coefficient-extraction formula or a
triangular recurrence, and if fixed-order remainders are controlled.
\end{definition}

\begin{definition}[Borel or hyperasymptotic completion]
A Borel completion supplements the divergent algebraic series by the Borel
singularities, terminant contributions, and exponentially small optimal
remainders that occur beyond all algebraic orders.
\end{definition}

The main formulas of this article are algebraically complete in the first sense.
For the forward correction we also give its exact Borel transform and pole
lattice.  This determines the beyond-all-orders scale and transports it through
the inverse map.  We do not disguise a terminant expansion as an ordinary
power series: the algebraic, iterated-logarithmic, and hyperasymptotic levels are
kept separate.

\begin{resultbox}
For both parities, all algebraic inverse coefficients are generated by one
universal recurrence.  The only parity dependence is the replacement
\[
 \log y-\tfrac12\log\pi
 \quad\longleftrightarrow\quad
 \log y-\tfrac12\log2
\]
inside the Lambert-\(W\) core.
\end{resultbox}

\chapter{Parity-resolved analytic models}

\section{Gamma representations}

For \(m\ge0\),
\[
 (2m)!!=2^m m!,
 \qquad
 (2m+1)!!=\frac{2^{m+1}\Gamma(m+\tfrac32)}{\sqrt\pi}.
\]
Equivalently, with a parity label \(\eps\in\{0,1\}\), define
\begin{equation}\label{eq:D-eps}
 D_\eps(x)
 :=C_\eps 2^{x/2}\Gamma\!\left(\frac{x}{2}+1\right),
 \qquad
 C_\eps:=\left(\frac{2}{\pi}\right)^{\eps/2}.
\end{equation}
Then
\[
 D_0(2m)=(2m)!!,
 \qquad
 D_1(2m+1)=(2m+1)!!.
\]
The two branch functions differ only by a constant:
\[
 D_1(x)=\sqrt{\frac{2}{\pi}}\,D_0(x).
\]

\begin{proposition}[Monotonicity]\label{prop:monotone}
Each \(D_\eps\) is strictly increasing on \([0,\infty)\), hence has a real
inverse on its range.
\end{proposition}

\begin{proof}
Logarithmic differentiation gives
\[
 \frac{D_\eps'(x)}{D_\eps(x)}
 =\frac12\left(
   \log2+\psi\!\left(\frac{x}{2}+1\right)
 \right),
\]
where \(\psi=\Gamma'/\Gamma\).  The trigamma function is positive on the
positive real axis, so \(\psi\) is increasing.  At \(x=0\), the bracket is
\(\log2-\gamma>0\).  It remains positive for every \(x\ge0\).
\end{proof}

\section{The centered coordinate}

The ordinary large-\(x\) Stirling expansion contains a half-logarithm:
\[
 \log D_\eps(x)
 \sim
 \frac{x}{2}(\log x-1)+\frac12\log x+A_\eps+\cdots.
\]
For inversion it is advantageous to absorb that half-logarithm into the
power--logarithmic phase.  Set
\begin{equation}\label{eq:u-centered}
 u:=x+1,
 \qquad
 A_\eps:=\log(C_\eps\sqrt\pi)
 =\begin{cases}
   \frac12\log\pi,&\eps=0,\\[2mm]
   \frac12\log2,&\eps=1.
  \end{cases}
\end{equation}
The choice \(u=x+1\) is not cosmetic: it converts the gamma argument into
\((u+1)/2=u/2+1/2\), and Bernoulli polynomials at \(1/2\) kill every odd
correction.

\begin{theorem}[Exact centered normal form]\label{thm:exact-normal}
For \(x>-1\), \(u=x+1\), and \(\eps\in\{0,1\}\),
\begin{equation}\label{eq:exact-factorization}
 D_\eps(x)
 =\exp(A_\eps)
  \left(\frac{u}{\e}\right)^{u/2}
  \exp\!\bigl(\Omega(u)\bigr),
\end{equation}
where
\begin{equation}\label{eq:Omega-def}
 \Omega(u):=
 \log\Gamma\!\left(\frac{u+1}{2}\right)
 -\frac{u}{2}\log\!\left(\frac{u}{2}\right)
 +\frac{u}{2}-\frac12\log(2\pi).
\end{equation}
In particular, \(\Omega\) is independent of parity.
\end{theorem}

\begin{proof}
Insert \(x=u-1\) in \cref{eq:D-eps}:
\[
 \log D_\eps(x)
 =\log C_\eps+\frac{u-1}{2}\log2
  +\log\Gamma\!\left(\frac{u+1}{2}\right).
\]
Subtract \(A_\eps+\tfrac{u}{2}(\log u-1)\), use
\(A_\eps=\log C_\eps+\tfrac12\log\pi\), and simplify.  The remainder is exactly
\(\Omega(u)\) in \cref{eq:Omega-def}.
\end{proof}

\section{A single parity transmonomial}

For integer \(n\), the constants in \cref{eq:u-centered} can be encoded by
\[
 A(n)=\frac14\log(2\pi)
      +\frac{(-1)^n}{4}\log\!\left(\frac{\pi}{2}\right).
\]
Consequently the exact sequence-level factorization is
\begin{equation}\label{eq:sequence-exact}
 n!!
 =\left(\frac{n+1}{\e}\right)^{(n+1)/2}
  \exp\!\left\{
   \frac14\log(2\pi)
   +\frac{(-1)^n}{4}\log\!\left(\frac{\pi}{2}\right)
   +\Omega(n+1)
  \right\}.
\end{equation}
Thus \((-1)^n\) is a genuine parity transmonomial.  The algebraic correction
\(\Omega(n+1)\) is common to both sectors.

\chapter{Exact Borel structure of the centered correction}

\section{A Binet-type integral}

\begin{theorem}[Borel--Laplace representation]\label{thm:borel}
For \(\Re u>0\),
\begin{equation}\label{eq:borel-laplace}
 \Omega(u)
 =\int_0^\infty \e^{-us}\,\widehat\Omega(s)\,ds,
 \qquad
 \widehat\Omega(s)
 :=\frac{1}{2s^2}\left(\frac{s}{\sinh s}-1\right).
\end{equation}
The apparent singularity of \(\widehat\Omega\) at \(s=0\) is removable.
\end{theorem}

\begin{proof}
It is convenient to put \(t=u/2\) and write
\[
 \widetilde\Omega(t)
 :=\log\Gamma(t+\tfrac12)-t\log t+t-\tfrac12\log(2\pi).
\]
The standard integral representation of the digamma function, combined with
Frullani's formula for \(\log t\), gives
\begin{align*}
 \psi(t+\tfrac12)-\log t
 &=\int_0^\infty \e^{-tr}
 \left(\frac1r-\frac{1}{2\sinh(r/2)}\right)\,dr.
\end{align*}
On the other hand,
\[
 \frac{d}{dt}
 \int_0^\infty \e^{-tr}
 \frac{r/(2\sinh(r/2))-1}{r^2}\,dr
 =
 \int_0^\infty \e^{-tr}
 \left(\frac1r-\frac{1}{2\sinh(r/2)}\right)\,dr.
\]
Both this integral and \(\widetilde\Omega(t)\) tend to zero as
\(t\to+\infty\); hence they are equal.  Finally substitute \(r=2s\).
The Taylor expansion \(s/\sinh s=1-s^2/6+\Ord(s^4)\) removes the origin.
\end{proof}

\section{Mittag--Leffler decomposition}

The elementary Mittag--Leffler expansion
\begin{equation}\label{eq:mittag}
 \frac{s}{\sinh s}
 =1+2s^2\sum_{m=1}^\infty
 \frac{(-1)^m}{s^2+\pi^2m^2}
\end{equation}
immediately gives
\begin{equation}\label{eq:borel-partial-fraction}
 \widehat\Omega(s)
 =\sum_{m=1}^\infty
 \frac{(-1)^m}{s^2+\pi^2m^2}.
\end{equation}
The Borel singularities are therefore the simple poles
\[
 s=\pm\ii\pi m,\qquad m=1,2,\ldots.
\]
Their residues are
\begin{align}
 \Res_{s=\ii\pi m}\widehat\Omega(s)
 &=\frac{(-1)^{m+1}\ii}{2\pi m},\label{eq:res-plus}\\
 \Res_{s=-\ii\pi m}\widehat\Omega(s)
 &=\frac{(-1)^m\ii}{2\pi m}.\label{eq:res-minus}
\end{align}
The nearest pair, at distance \(\pi\) from the origin, controls the
factorial divergence and the optimal remainder scale.

\chapter{The forward double-factorial transseries}

\section{The logarithmic coefficients}

Let \(B_n\) denote the Bernoulli numbers, with \(B_1=-1/2\), and let
\[
 \etaD(s)=\sum_{m=1}^\infty\frac{(-1)^{m-1}}{m^s}
 =(1-2^{1-s})\zeta(s)
\]
be the Dirichlet eta function.

\begin{theorem}[All-orders logarithmic expansion]\label{thm:forward-log}
As \(u\to\infty\) in any closed subsector of \(|\arg u|<\pi/2\),
\begin{equation}\label{eq:Omega-series}
 \Omega(u)\sim\sum_{k=1}^\infty\frac{c_k}{u^{2k-1}},
\end{equation}
where the following three formulas are equivalent:
\begin{equation}\label{eq:ck-three}
 \boxed{
 c_k=
 \frac{2^{2k-1}B_{2k}(\tfrac12)}{2k(2k-1)}
 =
 \frac{(1-2^{2k-1})B_{2k}}{2k(2k-1)}
 =
 (-1)^k\frac{(2k-2)!\etaD(2k)}{\pi^{2k}}.}
\end{equation}
For real \(u>0\), truncation after \(N\) terms has the exact remainder
\begin{equation}\label{eq:RN-exact}
 R_N(u)
 =(-1)^N\int_0^\infty \e^{-us}s^{2N}
 \sum_{m=1}^\infty
 \frac{(-1)^m}{(\pi m)^{2N}(s^2+\pi^2m^2)}\,ds,
\end{equation}
and therefore
\begin{equation}\label{eq:RN-sign-bound}
 0<(-1)^{N+1}R_N(u)
 <\frac{(2N)!}{\pi^{2N+2}u^{2N+1}}.
\end{equation}
\end{theorem}

\begin{proof}
Expand each summand of \cref{eq:borel-partial-fraction}:
\[
 \frac1{s^2+a^2}
 =\sum_{j=0}^{N-1}\frac{(-1)^js^{2j}}{a^{2j+2}}
 +\frac{(-1)^Ns^{2N}}{a^{2N}(s^2+a^2)}.
\]
Termwise Laplace integration of the finite sum gives
\[
 c_{j+1}
 =(2j)!\frac{(-1)^{j+1}\etaD(2j+2)}{\pi^{2j+2}},
\]
which is the last expression in \cref{eq:ck-three}.  Euler's formula for
\(B_{2k}\) and
\(B_{2k}(1/2)=(2^{1-2k}-1)B_{2k}\) give the other two forms
\cite{DLMF-Bernoulli}.  The remaining geometric-series term is
\cref{eq:RN-exact}.

For fixed \(s>0\), the positive sequence
\[
 \frac{1}{m^{2N}(s^2+\pi^2m^2)}
\]
strictly decreases with \(m\).  Its alternating sum, beginning with a negative
term, is negative and has magnitude smaller than its first term.  This proves
the sign assertion.  Finally,
\[
 |R_N(u)|
 <\int_0^\infty
 \frac{\e^{-us}s^{2N}}{\pi^{2N}(s^2+\pi^2)}\,ds
 <\frac1{\pi^{2N+2}}
 \int_0^\infty\e^{-us}s^{2N}\,ds,
\]
which is \cref{eq:RN-sign-bound}.
\end{proof}

The first coefficients are
\begin{equation}\label{eq:c-first}
 \Omega(u)\sim
 -\frac{1}{12u}
 +\frac{7}{360u^3}
 -\frac{31}{1260u^5}
 +\frac{127}{1680u^7}
 -\frac{511}{1188u^9}
 +\frac{1414477}{360360u^{11}}
 -\cdots.
\end{equation}

\begin{corollary}[Forward logarithmic transseries]\label{cor:forward-log}
For \(x\to+\infty\), \(u=x+1\), and fixed \(\eps\),
\begin{equation}\label{eq:logD-final}
 \boxed{
 \log D_\eps(x)
 \sim
 \frac{u}{2}(\log u-1)+A_\eps
 +\sum_{k=1}^\infty\frac{c_k}{u^{2k-1}}.}
\end{equation}
For integer \(n\), replace \(A_\eps\) by \(A(n)\) from
\cref{eq:sequence-exact}.
\end{corollary}

\section{Late coefficients and optimal truncation}

From \cref{eq:ck-three},
\begin{equation}\label{eq:ck-late}
 c_k
 =(-1)^k\frac{(2k-2)!}{\pi^{2k}}
 \left(1+\Ord(2^{-2k})\right).
\end{equation}
Thus the logarithmic series is Gevrey--1 when indexed by its actual power.
The ratio of consecutive term magnitudes is asymptotic to
\[
 \frac{(2k)(2k-1)}{\pi^2u^2}.
\]
The least term occurs near \(2k\approx\pi u\), and its size is
\begin{equation}\label{eq:least-term}
 \asymp u^{-1/2}\e^{-\pi u}.
\end{equation}
This agrees with the distance \(\pi\) to the nearest Borel poles.  On the
positive real axis the series alternates and is Borel summable without a
lateral ambiguity.  In complex directions, the poles
\(\pm\ii\pi m\) generate the usual terminant and Stokes structure for shifted
gamma asymptotics; see \cite{Nemes2015,ParisKaminski,Olver}.

\section{Multiplicative coefficients by Bell polynomials}

Define
\begin{equation}\label{eq:gamma-j}
 \gamma_j:=
 \begin{cases}
  c_{(j+1)/2},&j\text{ odd},\\
  0,&j\text{ even}.
 \end{cases}
\end{equation}
Then
\[
 \exp\!\left(\sum_{k\ge1}c_ku^{-(2k-1)}\right)
 =\sum_{m\ge0}\frac{s_m}{u^m}.
\]

\begin{theorem}[Complete Bell-polynomial formula]\label{thm:forward-bell}
The multiplicative coefficients are given by any of the equivalent formulas
\begin{align}
 s_m
 &=\frac1{m!}
   Y_m(1!\gamma_1,2!\gamma_2,\ldots,m!\gamma_m),
   \label{eq:sm-complete-bell}\\
 &=\sum_{\substack{k_1,\ldots,k_m\ge0\\
                   k_1+2k_2+\cdots+mk_m=m}}
   \prod_{j=1}^m\frac{\gamma_j^{k_j}}{k_j!},
   \label{eq:sm-partition}\\
 s_0&=1,\qquad
 m s_m=\sum_{j=1}^m j\gamma_j s_{m-j}
 \quad(m\ge1).
 \label{eq:sm-recurrence}
\end{align}
Consequently,
\begin{equation}\label{eq:forward-mult}
 \boxed{
 D_\eps(x)
 \sim
 \e^{A_\eps}
 \left(\frac{u}{\e}\right)^{u/2}
 \sum_{m=0}^\infty\frac{s_m}{u^m}.}
\end{equation}
\end{theorem}

\begin{proof}
Put \(t=u^{-1}\) and
\(S(t)=\exp(\sum_{j\ge1}\gamma_jt^j)\).  The exponential formula for complete
Bell polynomials gives \cref{eq:sm-complete-bell}; direct multinomial expansion
gives \cref{eq:sm-partition}.  Differentiating
\(S(t)\) and comparing coefficients in
\[
 tS'(t)=
 \left(\sum_{j\ge1}j\gamma_jt^j\right)S(t)
\]
gives \cref{eq:sm-recurrence}.
\end{proof}

The first coefficients are displayed in \cref{tab:s-coeff}.
\begin{table}[ht]
\centering
\caption{First multiplicative coefficients in \cref{eq:forward-mult}.}
\label{tab:s-coeff}
\begin{tabular}{c@{\qquad}l}
\toprule
\(m\)&\(s_m\)\\
\midrule
0&\(1\)\\
1&\(-1/12\)\\
2&\(1/288\)\\
3&\(1003/51840\)\\
4&\(-4027/2488320\)\\
5&\(-5128423/209018880\)\\
6&\(168359651/75246796800\)\\
7&\(68168266699/902961561600\)\\
8&\(-587283555451/86684309913600\)\\
\bottomrule
\end{tabular}
\end{table}

\section{The uncentered form}

For comparison with conventional Stirling formulas, direct expansion at
\(x=\infty\) gives
\begin{equation}\label{eq:uncentered-log}
 \log D_\eps(x)
 \sim
 \frac{x}{2}(\log x-1)+\frac12\log x+A_\eps
 +\sum_{k=1}^\infty
 \frac{b_k}{x^{2k-1}},
\end{equation}
where
\begin{equation}\label{eq:bk}
 b_k=\frac{2^{2k-1}B_{2k}}{2k(2k-1)}.
\end{equation}
The leading term is therefore
\[
 D_\eps(x)
 \sim \e^{A_\eps}\sqrt{x}\left(\frac{x}{\e}\right)^{x/2},
\]
which yields the familiar constants \(\sqrt\pi\) for even \(n\) and
\(\sqrt2\) for odd \(n\) \cite{OEIS-A006882}.

The centered and uncentered coefficients satisfy the identity
\begin{equation}\label{eq:center-shift-identity}
 \frac{(-1)^{r+1}}{2r(r+1)}
 +(-1)^{r+1}
 \sum_{k=1}^{\lfloor(r+1)/2\rfloor}
 \binom{r-1}{2k-2}c_k
 =
 \begin{cases}
  b_{(r+1)/2},&r\text{ odd},\\
  0,&r\text{ even}.
 \end{cases}
\end{equation}
It follows simply by substituting \(u=x+1\) into
\cref{eq:logD-final} and comparing powers.  The cancellation of every even
power is another way to see why the shift \(u=x+1\) is natural.

\chapter{The exact Lambert core}

\section{Dominant phase inversion}

Fix \(\eps\in\{0,1\}\) and let \(y\to+\infty\).  Set
\begin{equation}\label{eq:R-eps}
 R_\eps:=\log y-A_\eps.
\end{equation}
Ignoring \(\Omega(u)\) in \cref{eq:exact-factorization} leaves
\begin{equation}\label{eq:core-equation}
 \frac{u}{2}(\log u-1)=R_\eps.
\end{equation}

\begin{proposition}[Lambert solution of the core]\label{prop:Lambert-core}
Define
\begin{equation}\label{eq:v-X-Lambda}
 v:=\W_0\!\left(\frac{2R_\eps}{\e}\right),
 \qquad
 X:=\frac{2R_\eps}{v}=\e^{v+1},
 \qquad
 \Lambda:=v+1=\log X.
\end{equation}
Then \(X\) is the unique positive solution of \cref{eq:core-equation}, and
\[
 \frac{d}{du}\left[\frac{u}{2}(\log u-1)\right]_{u=X}
 =\frac{\Lambda}{2}.
\]
\end{proposition}

\begin{proof}
Writing \(v=\log u-1\) gives \(u=\e^{v+1}\), and
\cref{eq:core-equation} becomes
\[
 v\e^v=\frac{2R_\eps}{\e}.
\]
For positive \(R_\eps\) the principal branch \(\W_0\) is real and positive.
The derivative formula is immediate.
\end{proof}

The exact Lambert core avoids a premature expansion in nested logarithms.
This is both algebraically cleaner and numerically more stable.  Iterated
logarithms will be introduced only after the inverse-power correction blocks
have been generated.

\section{Scale hierarchy}

As \(y\to\infty\),
\[
 R_\eps\sim\log y,\qquad
 v\sim\log R_\eps,\qquad
 X\sim\frac{2R_\eps}{\log R_\eps},\qquad
 \Lambda\sim\log R_\eps.
\]
The first forward correction is \(\Omega(X)\sim-1/(12X)\).  Division by the
core slope \(\Lambda/2\) shows that the first displacement of the inverse is
of size
\[
 \frac{1}{6\Lambda X}\sim\frac{1}{12R_\eps}.
\]
This scale is invisible if one retains only the leading Lambert core.

\chapter{Universal all-orders inversion}

\section{The normalized implicit equation}

Write the exact solution in relative form
\begin{equation}\label{eq:E-def}
 u=X(1+E),\qquad q:=X^{-1}.
\end{equation}
Define
\begin{equation}\label{eq:H-def}
 H(E):=(1+E)\log(1+E)-E
 =\sum_{m=2}^\infty\frac{(-1)^m}{m(m-1)}E^m
\end{equation}
and
\begin{equation}\label{eq:ak}
 a_k:=2c_k
 =\frac{(1-2^{2k-1})B_{2k}}{k(2k-1)}.
\end{equation}

\begin{theorem}[Master inverse equation]\label{thm:inverse-master}
After formal substitution of the complete asymptotic series for \(\Omega\),
the inverse equation is
\begin{equation}\label{eq:master-normalized}
 \boxed{
 \Lambda E+H(E)
 +\sum_{k=1}^\infty
 a_kq^{2k}(1+E)^{-(2k-1)}=0.}
\end{equation}
It has a unique formal solution of the form
\begin{equation}\label{eq:E-series}
 E(q;\Lambda)=\sum_{n=1}^\infty e_n(\Lambda)q^{2n}.
\end{equation}
Therefore
\begin{equation}\label{eq:inverse-series}
 \boxed{
 D_\eps^{-1}(y)
 \sim X-1+\sum_{n=1}^\infty
 \frac{d_{2n-1}(\Lambda)}{X^{2n-1}},
 \qquad d_{2n-1}=e_n.}
\end{equation}
No even inverse powers of \(X\) occur.
\end{theorem}

\begin{proof}
Because \(X\) solves the core equation,
\begin{align*}
 &\frac{X(1+E)}{2}
 \bigl(\log(X(1+E))-1\bigr)
 -\frac{X}{2}(\log X-1)\\
 &\qquad
 =\frac{X}{2}\left[
   \Lambda E+(1+E)\log(1+E)-E
  \right]
 =\frac{X}{2}\bigl(\Lambda E+H(E)\bigr).
\end{align*}
The exact equation is the core difference plus
\(\Omega(X(1+E))=0\).  Multiplying by \(2/X=2q\), then substituting
\[
 \Omega(X(1+E))
 \sim\sum_{k\ge1}c_kq^{2k-1}(1+E)^{-(2k-1)}
\]
gives \cref{eq:master-normalized}.  Its linear term is \(\Lambda E\),
with \(\Lambda\ne0\).  Formal implicit inversion therefore gives a unique
solution.  Since every forcing term begins with an even power of \(q\), and
\(H\) has no constant or linear term, induction shows that \(E\) contains only
\(q^{2n}\).  Finally, \(u-X=XE=\sum e_nX^{1-2n}\), and \(x=u-1\).
\end{proof}

\section{Triangular coefficient extraction}

Let
\[
 E_{n-1}(q):=\sum_{j=1}^{n-1}e_j(\Lambda)q^{2j}.
\]

\begin{theorem}[Universal triangular recurrence]\label{thm:inverse-recursion}
For \(n\ge1\),
\begin{equation}\label{eq:e-recursion}
 \boxed{
 e_n(\Lambda)=
 -\frac1\Lambda\coeff{q^{2n}}{
  H(E_{n-1})
  +\sum_{k=1}^n
   a_kq^{2k}(1+E_{n-1})^{-(2k-1)}
 }.}
\end{equation}
Every coefficient on the right is already known when \(e_n\) is computed.
\end{theorem}

\begin{proof}
Insert
\(E=E_{n-1}+e_nq^{2n}+\Ord(q^{2n+2})\) into
\cref{eq:master-normalized}.  The term \(\Lambda E\) contributes
\(\Lambda e_nq^{2n}\).  Since \(H(E)=\Ord(E^2)\), no occurrence of
\(e_nq^{2n}\) can contribute to the coefficient of \(q^{2n}\) in \(H(E)\).
In the \(k\)-th forcing term the required residual degree is \(2(n-k)<2n\),
so \(e_n\) again cannot occur.  Solving the coefficient equation gives
\cref{eq:e-recursion}.
\end{proof}

\begin{methodbox}
The recurrence is finite at every order.  To obtain \(e_n\), one needs only
\(a_1,\ldots,a_n\) and \(e_1,\ldots,e_{n-1}\).  No asymptotic guessing,
numerical fitting, or re-expansion of Lambert \(W\) is required.
\end{methodbox}

\section{Scalar Bell-polynomial formula}

For a series \(E(z)=\sum_{j\ge1}e_jz^j\), define the ordinary partial Bell
convolution
\begin{equation}\label{eq:ordinary-Bell}
 \widehat B_{r,m}(e_1,e_2,\ldots)
 :=\coeff{z^r}{E(z)^m}
 =\frac{m!}{r!}
 B_{r,m}(1!e_1,2!e_2,\ldots),
\end{equation}
with
\[
 \widehat B_{0,0}=1,\qquad
 \widehat B_{r,0}=0\ (r>0).
\]

\begin{corollary}[Bell-polynomial recurrence]\label{cor:Bell-recurrence}
The coefficients in \cref{eq:E-series} satisfy
\begin{align}
 e_n=-\frac1\Lambda\Bigg[
 &\sum_{m=2}^n
 \frac{(-1)^m}{m(m-1)}
 \widehat B_{n,m}\notag\\
 &+\sum_{k=1}^n a_k
 \sum_{m=0}^{n-k}
 (-1)^m\binom{2k+m-2}{m}
 \widehat B_{n-k,m}
 \Bigg].
 \label{eq:e-Bell}
\end{align}
All Bell polynomials in this formula involve only
\(e_1,\ldots,e_{n-1}\).
\end{corollary}

\begin{proof}
Use the series for \(H\) in \cref{eq:H-def} and
\[
 (1+E)^{-(2k-1)}
 =\sum_{m\ge0}
 (-1)^m\binom{2k+m-2}{m}E^m.
\]
Extract the coefficient of \(z^n\), where \(z=q^2\), and apply
\cref{eq:ordinary-Bell}.
\end{proof}

\section{Denominator and leading-term structure}

\begin{theorem}[Polynomial structure]\label{thm:polynomial-structure}
For every \(n\ge1\),
\begin{equation}\label{eq:en-ring}
 e_n(\Lambda)\in
 \Lambda^{-1}\Q[\Lambda^{-1}],
 \qquad
 \deg_{\Lambda^{-1}}e_n\le2n-1.
\end{equation}
More precisely,
\begin{equation}\label{eq:en-leading}
 e_n(\Lambda)
 =-\frac{a_n}{\Lambda}
  +\Ord(\Lambda^{-2})
 \qquad(\Lambda\to\infty).
\end{equation}
Equivalently,
\[
 d_{2n-1}(\Lambda)
 =-\frac{2c_n}{\Lambda}+\Ord(\Lambda^{-2}).
\]
\end{theorem}

\begin{proof}
The assertion is immediate for \(n=1\).  Assume it for all smaller indices.
Every monomial in \(\widehat B_{r,m}\) is a product of \(m\) earlier
coefficients whose indices sum to \(r\).  Its least power of
\(\Lambda^{-1}\) is \(m\), while its largest is at most
\[
 \sum_{\nu=1}^m(2j_\nu-1)=2r-m.
\]
The outer factor \(\Lambda^{-1}\) in \cref{eq:e-Bell} raises the largest
possible exponent by one.  In the first sum, \(r=n\) and \(m\ge2\), giving at
most \(2n-1\).  In the forcing sum, \(r=n-k\), and multiplication by \(a_k\)
does not affect \(\Lambda\); the same bound follows.  Rationality is preserved
by every operation.

For the leading term as \(\Lambda\to\infty\), all Bell terms contain at least
one earlier \(e_j=\Ord(\Lambda^{-1})\).  The only term of order
\(\Lambda^{-1}\) after the outer division is the \(k=n,m=0\) forcing term,
namely \(-a_n/\Lambda\).
\end{proof}

\section{A nonrecursive Lagrange--B\"urmann formula}

Define the Bernoulli forcing series
\begin{equation}\label{eq:A-series}
 A(z):=\sum_{k=1}^\infty a_kz^k,
 \qquad a_1=-\frac16,
\end{equation}
and let \(S=A^{\langle-1\rangle}\) be its formal compositional inverse.

Equation \cref{eq:master-normalized}, with \(z=q^2\), can be rewritten as
\begin{equation}\label{eq:A-compressed}
 \Lambda E+H(E)
 +(1+E)A\!\left(\frac{z}{(1+E)^2}\right)=0.
\end{equation}
Consequently,
\begin{equation}\label{eq:G-Lambda}
 z=G_\Lambda(E):=
 (1+E)^2
 S\!\left(
  -\frac{\Lambda E+H(E)}{1+E}
 \right).
\end{equation}
Since \(G_\Lambda'(0)=-\Lambda/a_1=6\Lambda\ne0\), this series is revertible.

\begin{theorem}[Closed coefficient extraction]\label{thm:lagrange-closed}
For every \(n\ge1\),
\begin{equation}\label{eq:e-lagrange}
 \boxed{
 e_n(\Lambda)
 =\frac1n\coeff{E^{n-1}}{
  \left(\frac{E}{G_\Lambda(E)}\right)^n}.}
\end{equation}
The auxiliary inverse \(S\) itself has coefficients
\begin{equation}\label{eq:S-lagrange}
 \coeff{w^m}{S(w)}
 =\frac1m\coeff{t^{m-1}}{
   \left(\frac{t}{A(t)}\right)^m}.
\end{equation}
Thus \cref{eq:e-lagrange,eq:S-lagrange} form a nonrecursive finite
coefficient formula for every \(e_n\).
\end{theorem}

\begin{proof}
Apply the formal Lagrange--B\"urmann theorem first to \(z=G_\Lambda(E)\) and
then to \(w=A(t)\).  Although \(A\) is a divergent Gevrey series, it is a valid
formal power series with nonzero linear coefficient.  Every requested finite
jet uses only finitely many \(a_k\), so no analytic convergence assumption is
needed.
\end{proof}

\chapter{Explicit inverse blocks}

\section{The first five universal coefficients}

Repeated use of \cref{eq:e-recursion} gives
\begin{align}
 d_1(\Lambda)
 &=\frac{1}{6\Lambda},\label{eq:d1}\\[1mm]
 d_3(\Lambda)
 &=-\frac{14\Lambda^2+10\Lambda+5}
 {360\Lambda^3},\label{eq:d3}\\[1mm]
 d_5(\Lambda)
 &=\frac{
 2232\Lambda^4+1176\Lambda^3+714\Lambda^2
 +350\Lambda+105}
 {45360\Lambda^5},\label{eq:d5}\\[1mm]
 d_7(\Lambda)
 &=-\frac{1}{5443200\Lambda^7}
 \Bigl(
 822960\Lambda^6+292536\Lambda^5+136956\Lambda^4
 \notag\\[-1mm]
 &\hspace{28mm}
 +68040\Lambda^3+32970\Lambda^2+12250\Lambda+2625
 \Bigr),\label{eq:d7}\\[1mm]
 d_9(\Lambda)
 &=\frac{1}{359251200\Lambda^9}
 \Bigl(
 309052800\Lambda^8+77920128\Lambda^7
 +26024592\Lambda^6
 \notag\\[-1mm]
 &\hspace{15mm}
 +10019592\Lambda^5+4516028\Lambda^4
 +2023560\Lambda^3
 \notag\\[-1mm]
 &\hspace{15mm}
 +812350\Lambda^2+242550\Lambda+40425
 \Bigr).
 \label{eq:d9}
\end{align}
Combining these blocks gives the compact expansion through \(X^{-9}\):
\begin{equation}\label{eq:inverse-explicit-clean}
 D_\eps^{-1}(y)
 \sim X-1+\frac{d_1(\Lambda)}{X}
 +\frac{d_3(\Lambda)}{X^3}
 +\frac{d_5(\Lambda)}{X^5}
 +\frac{d_7(\Lambda)}{X^7}
 +\frac{d_9(\Lambda)}{X^9}+\cdots.
\end{equation}

\section{A compact \texorpdfstring{Lambert-\(W\)}{Lambert-W} version}

Because \(X=2R_\eps/v\) and \(\Lambda=v+1\), the first terms can be written
without \(X\):
\begin{align}
 D_\eps^{-1}(y)
 \sim{}&
 \frac{2R_\eps}{v}-1
 +\frac{v}{12R_\eps(v+1)}
 \notag\\
 &-\frac{
 v^3\left(14(v+1)^2+10(v+1)+5\right)}
 {2880R_\eps^3(v+1)^3}
 \notag\\
 &+\frac{
 v^5\left(
 2232(v+1)^4+1176(v+1)^3+714(v+1)^2
 +350(v+1)+105
 \right)}
 {1451520R_\eps^5(v+1)^5}
 +\cdots,
 \label{eq:inverse-v}
\end{align}
where \(v=\W_0(2R_\eps/\e)\).  Keeping \(v\) unexpanded is generally the best
numerical form.

\section{A correction to a common first guess}

A frequently used rough inversion is
\[
 x\approx\frac{2R_\eps}{\W_0(2R_\eps/\e)}-1.
\]
The leading correction is not of order \(1/X\) alone but
\[
 \frac{1}{6\Lambda X}
 =\frac{v}{12R_\eps(v+1)}
 =\frac{1}{12R_\eps}
  \left(1-\frac{1}{v+1}\right).
\]
This distinction matters when numerical inversion is desired: the factor
\(\Lambda^{-1}\) is forced by the derivative of the power--logarithmic phase.

\chapter{Fixed-order error control}

\section{Residual-to-root conversion}

Let
\[
 u_N:=X+\sum_{n=1}^N e_n(\Lambda)X^{1-2n},
 \qquad x_N:=u_N-1.
\]

\begin{theorem}[Poincar\'e validity of the inverse expansion]
\label{thm:inverse-remainder}
For every fixed \(N\ge0\),
\begin{equation}\label{eq:inverse-remainder}
 D_\eps^{-1}(y)
 =X-1+\sum_{n=1}^N
  e_n(\Lambda)X^{1-2n}
 +\Ord\!\left(
   \frac{X^{-(2N+1)}}{\Lambda}
  \right)
\end{equation}
as \(y\to+\infty\).  The implied constant can be chosen uniformly for
\(\eps=0,1\).
\end{theorem}

\begin{proof}
By construction, substitution of \(E_N=(u_N-X)/X\) into the normalized
equation cancels every term through \(q^{2N}\).  The first uncancelled term is
\(\Ord(q^{2N+2})\), with coefficients that are rational functions of
\(\Lambda\) bounded for large \(\Lambda\).  Undoing the normalization
multiplies the phase residual by \(X/2\), hence the residual in
\(\log D_\eps\) is \(\Ord(X^{-(2N+1)})\).

For \(u\) between \(u_N\) and the true root,
\[
 \frac{d}{du}
 \left[\frac{u}{2}(\log u-1)+\Omega(u)\right]
 =\frac12\log u+\Omega'(u)
 =\frac12\Lambda\,(1+\littleo(1)).
\]
The mean-value theorem therefore divides the phase residual by a quantity
asymptotic to \(\Lambda/2\), proving \cref{eq:inverse-remainder}.  The
parity constant disappears upon differentiation, so the estimate is uniform
in \(\eps\).
\end{proof}

\section{Backward error as a certification device}

For numerical work, a truncated approximation \(\widetilde x\) should be
certified by its logarithmic residual
\[
 \rho_\eps(\widetilde x;y)
 :=\log D_\eps(\widetilde x)-\log y.
\]
Since the branch is increasing,
\begin{equation}\label{eq:backward-to-forward}
 |\widetilde x-D_\eps^{-1}(y)|
 \le
 \frac{2|\rho_\eps(\widetilde x;y)|}
 {\inf_{\xi}\{\log2+\psi(\xi/2+1)\}},
\end{equation}
where the infimum is over the interval between the approximation and the
root.  At large arguments the denominator is \(\sim\Lambda\), matching
\cref{eq:inverse-remainder}.  This backward-error test is more reliable than
estimating the forward error from the last retained term alone.

\chapter{The complete iterated-logarithm expansion}

\section{\texorpdfstring{Lambert \(W\)}{Lambert W} in Stirling-cycle form}

Let
\begin{equation}\label{eq:L1L2}
 Z:=\frac{2R_\eps}{\e},
 \qquad L_1:=\log Z,
 \qquad L_2:=\log L_1.
\end{equation}
The large-\(Z\) expansion of the principal Lambert function can be written
using unsigned Stirling cycle numbers
\(\genfrac{[}{]}{0pt}{}{n}{k}\):
\begin{equation}\label{eq:W-Stirling}
 \W_0(Z)
 \sim
 L_1-L_2
 +\sum_{n=1}^\infty\frac1{L_1^n}
 \sum_{m=1}^n
 (-1)^{n+m}
 \genfrac{[}{]}{0pt}{}{n}{n-m+1}
 \frac{L_2^m}{m!}.
\end{equation}
This is the standard Stirling-number form of the Lambert expansion
\cite{Corless1996,DLMF-W}.

Write
\begin{equation}\label{eq:alpha-def}
 \frac{\W_0(Z)}{L_1}
 =1+\sum_{j=1}^\infty
 \frac{\alpha_j(L_2)}{L_1^j}.
\end{equation}
Then
\begin{align}
 \alpha_1(L_2)&=-L_2,\label{eq:alpha1}\\
 \alpha_j(L_2)
 &=
 \sum_{m=1}^{j-1}
 (-1)^{j-1+m}
 \genfrac{[}{]}{0pt}{}{j-1}{j-m}
 \frac{L_2^m}{m!},
 \qquad j\ge2.
 \label{eq:alpha-general}
\end{align}

\section{Reciprocal coefficients}

Define \(\beta_0=1\) and
\begin{equation}\label{eq:beta-recurrence}
 \beta_n(L_2)
 =-\sum_{j=1}^n
 \alpha_j(L_2)\beta_{n-j}(L_2).
\end{equation}
Then
\begin{equation}\label{eq:X-outer}
 X
 \sim\frac{2R_\eps}{L_1}
 \sum_{n=0}^\infty\frac{\beta_n(L_2)}{L_1^n}.
\end{equation}
The first polynomials are
\begin{align*}
 \beta_0&=1,\\
 \beta_1&=L_2,\\
 \beta_2&=L_2^2-L_2,\\
 \beta_3&=L_2^3-\frac52L_2^2+L_2,\\
 \beta_4&=L_2^4-\frac{13}{3}L_2^3+\frac92L_2^2-L_2,\\
 \beta_5&=L_2^5-\frac{77}{12}L_2^4
          +\frac{71}{6}L_2^3-7L_2^2+L_2.
\end{align*}

\section{All coefficients in the two-level transseries}

Introduce \(t=L_1^{-1}\) and the formal polynomial-logarithmic series
\begin{equation}\label{eq:Wcal-Lcal}
 \mathcal W(t;L_2):=
 1+\sum_{j\ge1}\alpha_j(L_2)t^j,
 \qquad
 \mathcal L(t;L_2):=\mathcal W(t;L_2)+t.
\end{equation}
Then
\[
 \frac{W}{L_1}=\mathcal W,
 \qquad
 \frac{\Lambda}{L_1}=\mathcal L,
 \qquad
 X=\frac{2R_\eps}{L_1}\mathcal W^{-1}.
\]
Write
\begin{equation}\label{eq:en-coeff-p}
 e_n(\Lambda)=
 \sum_{p=1}^{2n-1}e_{n,p}\Lambda^{-p},
 \qquad e_{n,p}\in\Q.
\end{equation}

\begin{theorem}[Complete outer coefficient formula]\label{thm:outer-complete}
The full algebraic inverse transseries is
\begin{align}
 D_\eps^{-1}(y)
 \sim{}&
 \frac{2R_\eps}{L_1}
 \sum_{r\ge0}\beta_r(L_2)L_1^{-r}-1
 \notag\\
 &+\sum_{n\ge1}
 \left(\frac{L_1}{2R_\eps}\right)^{2n-1}
 \sum_{r\ge1}
 D_{n,r}(L_2)L_1^{-r},
 \label{eq:full-outer}
\end{align}
where every polynomial \(D_{n,r}\) is given by the finite extraction
\begin{equation}\label{eq:Dnr}
 \boxed{
 D_{n,r}(L_2)
 =
 \coeff{t^r}{
  \mathcal W(t;L_2)^{2n-1}
  \sum_{p=1}^{2n-1}
  e_{n,p}t^p\mathcal L(t;L_2)^{-p}
 }.}
\end{equation}
Together with \cref{eq:e-recursion,eq:alpha-general,eq:beta-recurrence},
this gives a general formula for every coefficient of every inverse sector.
\end{theorem}

\begin{proof}
The core term follows from
\(X=(2R_\eps/L_1)\mathcal W^{-1}\), and
\cref{eq:beta-recurrence} is precisely reciprocal-series division.
For the \(n\)-th inverse block,
\begin{align*}
 e_n(\Lambda)X^{-(2n-1)}
 &=
 \left(\frac{L_1}{2R_\eps}\right)^{2n-1}
 \mathcal W^{2n-1}
 \sum_{p=1}^{2n-1}
 e_{n,p}L_1^{-p}\mathcal L^{-p}.
\end{align*}
Set \(t=L_1^{-1}\) and extract its coefficient.  At any fixed \(r\), only
finitely many \(\alpha_j\) and \(e_{n,p}\) are needed.
\end{proof}

\begin{corollary}[Leading term of each algebraic sector]
The leading term in the \(n\)-th inverse-power sector is
\begin{equation}\label{eq:sector-leading}
 -a_n
 \frac{L_1^{2n-2}}{(2R_\eps)^{2n-1}}.
\end{equation}
\end{corollary}

\begin{proof}
In \cref{eq:Dnr}, the first nonzero power is \(r=1\), arising from
\(e_{n,1}=-a_n\) and the constant terms of \(\mathcal W,\mathcal L\).
\end{proof}

\begin{summarybox}
The inverse contains two nested grids:
\[
 \frac{R_\eps}{L_1}
 \sum_{r\ge0}\frac{\Q[L_2]}{L_1^r}
 \quad\text{and}\quad
 \sum_{n\ge1}
 \frac{L_1^{2n-1}}{R_\eps^{2n-1}}
 \sum_{r\ge1}\frac{\Q[L_2]}{L_1^r}.
\]
Unsigned Stirling numbers generate the first grid; Bernoulli numbers generate
the sector labels \(n\); Bell and Lagrange coefficients generate the internal
polynomials.
\end{summarybox}

\chapter{Beyond all algebraic orders}

\section{Optimal forward remainder}

The exact Borel kernel shows that the formal log series has nearest singularities
at \(\pm\ii\pi\).  If it is truncated near its least term,
\(N\approx\pi u/2\), then
\[
 R_N(u)=\Ord\!\left(u^{-1/2}\e^{-\pi u}\right)
 \qquad(u\to+\infty).
\]
A more refined exponentially improved expansion resolves \(R_N\) into
terminant contributions associated with the pole pairs
\(\pm\ii\pi m\).  Such expansions are standard for the gamma function and its
reciprocal \cite{Nemes2015}; the shifted form here has the explicit Borel
transform \cref{eq:borel-laplace}, so the singularity locations and residues
are already known.

\section{Transport through the inverse map}

Suppose an algebraically truncated forward phase is written
\[
 \Phi(u)=\frac{u}{2}(\log u-1)+A_\eps+\Omega_{\rm alg}(u)
          +\mathcal R(u),
\]
where \(\mathcal R\) is an exponentially improved remainder.  Let
\(u_{\rm alg}\) solve the equation without \(\mathcal R\).  One Newton
linearization gives
\begin{equation}\label{eq:exp-transport}
 u-u_{\rm alg}
 \sim
 -\frac{\mathcal R(u_{\rm alg})}
  {\frac12\log u_{\rm alg}+\Omega_{\rm alg}'(u_{\rm alg})}
 \sim
 -\frac{2\mathcal R(X)}{\Lambda}.
\end{equation}
Thus the optimally truncated inverse remainder on the positive real axis has
the scale
\begin{equation}\label{eq:inverse-exp-scale}
 \Ord\!\left(
 \frac{X^{-1/2}\e^{-\pi X}}{\Lambda}
 \right).
\end{equation}

For a full terminant transseries, retain the exact remainder in the normalized
equation:
\begin{equation}\label{eq:master-with-rem}
 \Lambda E+H(E)
 +\sum_{k\ge1}a_kq^{2k}(1+E)^{-(2k-1)}
 +2q\,\mathcal R(X(1+E))=0.
\end{equation}
Coefficient extraction in the enlarged transmonomial semigroup then proceeds
by the same Bell/Fa\`a di Bruno operations as before.  At first order, every
forward terminant sector is multiplied by \(-2/\Lambda\); nonlinear products
generate higher sectors.

\begin{warningbox}
The complex exponentials associated with Borel poles are
\(\e^{\mp\ii\pi m u}\), together with direction-dependent terminants.
On the positive real axis it is the terminant-weighted combination, not a bare
oscillatory exponential, whose optimally truncated magnitude is
\(\e^{-\pi u}\).  Conflating these two descriptions gives incorrect Stokes
statements.
\end{warningbox}

\chapter{The inverse of the discrete sequence}\label{ch:discrete}

\section{Parity candidates}

For \(y>1\), let
\[
 x_0(y):=D_0^{-1}(y),
 \qquad
 x_1(y):=D_1^{-1}(y).
\]
The least even and odd indices whose double factorials exceed \(y\) are
\begin{align}
 E(y)&:=2\left\lceil\frac{x_0(y)}{2}\right\rceil,
 \label{eq:even-candidate}\\
 O(y)&:=2\left\lceil\frac{x_1(y)-1}{2}\right\rceil+1.
 \label{eq:odd-candidate}
\end{align}

\begin{theorem}[Exact generalized inverse]\label{thm:discrete-inverse}
Let
\[
 N(y):=\min\{n\in\N:n!!\ge y\}.
\]
For \(y>1\),
\begin{equation}\label{eq:N-exact}
 \boxed{N(y)=\min\{E(y),O(y)\}.}
\end{equation}
The real quantities \(x_0,x_1\) in \cref{eq:even-candidate,eq:odd-candidate}
may be evaluated by the transseries \cref{eq:inverse-series}.
\end{theorem}

\begin{proof}
On each parity class, \(n!!\) is strictly increasing, so
\cref{eq:even-candidate} is the least admissible even index and
\cref{eq:odd-candidate} the least admissible odd index.  The least admissible
index without a parity restriction is their minimum.
\end{proof}

\section{Rounding and certification}

An asymptotic approximation can be rounded safely only when its error interval
does not cross a parity lattice point.  A robust procedure is:

\begin{algorithm}[Certified discrete inversion]
\begin{enumerate}[label=\arabic*.]
\item Compute truncated approximations to \(x_0(y)\) and \(x_1(y)\).
\item Form the two neighboring even integers around \(x_0\) and the two
      neighboring odd integers around \(x_1\).
\item Evaluate \(\log D_\eps(n)\) at those candidates, using either
      \(\log\Gamma\) or the exact recurrence.
\item Select the least candidate satisfying \(\log D_\eps(n)\ge\log y\).
\end{enumerate}
\end{algorithm}

Only a constant number of exact or high-precision checks is required.
The asymptotic expansion supplies the location; the final comparisons remove
all rounding ambiguity.

\chapter{A single analytic interpolation and its oscillatory inverse}
\label{ch:interpolation}

\section{The all-integer interpolation}

OEIS records the analytic interpolation
\begin{equation}\label{eq:Luschny}
 \mathscr D(x)
 :=
 2^{x/2}
 \left(\frac{2}{\pi}\right)^{
   \frac12\sin^2(\pi x/2)}
 \Gamma\!\left(\frac{x}{2}+1\right),
\end{equation}
which agrees with \(n!!\) at every nonnegative integer
\cite{OEIS-A006882}.  In terms of \(u=x+1\),
\begin{equation}\label{eq:interp-phase}
 \log\mathscr D(x)
 =
 \frac{u}{2}(\log u-1)+A_0+\Omega(u)+P(x),
\end{equation}
where
\begin{equation}\label{eq:P-periodic}
 P(x):=
 \frac14(1-\cos\pi x)\log\!\left(\frac{2}{\pi}\right).
\end{equation}
At even integers \(P=0\); at odd integers
\(P=\tfrac12\log(2/\pi)\), exactly converting \(A_0\) to \(A_1\).
Since
\[
 \frac{d}{dx}\log\mathscr D(x)
 =\frac12\!\left(\log2+\psi\!\left(\frac{x}{2}+1\right)\right)+P'(x)
 \sim\frac12\log x,
\]
and \(P'\) is bounded, \(\mathscr D\) is strictly increasing on a real tail.
Throughout this chapter, \(\mathscr D^{-1}\) denotes that tail inverse as
\(y\to+\infty\), not a globally single-valued inverse near the origin.

\begin{warningbox}
The periodic phase \(P(x)=\Ord(1)\) displaces the inverse by
\(\Ord(\Lambda^{-1})\).  The Bernoulli correction displaces it only by
\(\Ord((X\Lambda)^{-1})\).  Therefore the inverse of the single interpolation
\(\mathscr D\) is \emph{not} obtained by simply forgetting parity in the
branchwise formulas.
\end{warningbox}

\section{Phase-coordinate Lagrange inversion}

Let
\[
 R_0:=\log y-A_0,
 \qquad
 h(R):=\frac{2R}{\W_0(2R/\e)}.
\]
Thus \(h\) is the inverse of \(u\mapsto\tfrac12u(\log u-1)\).  Define
\begin{equation}\label{eq:phase-A}
 \mathcal A(R):=
 P(h(R)-1)+\Omega(h(R)).
\end{equation}
With \(s=\tfrac12u(\log u-1)\), the interpolation equation becomes
\[
 R_0=s+\mathcal A(s).
\]

\begin{theorem}[All-orders phase-coordinate reversion]
\label{thm:phase-reversion}
Introduce a bookkeeping parameter \(\tau\) in
\(R_0=s+\tau\mathcal A(s)\).  Its formal inverse is
\begin{equation}\label{eq:phase-s}
 s=
 R_0+\sum_{m=1}^\infty
 \frac{(-\tau)^m}{m!}
 \left(\frac{d}{dR_0}\right)^{m-1}
 \mathcal A(R_0)^m.
\end{equation}
Setting \(\tau=1\) as an asymptotic expansion gives
\begin{equation}\label{eq:interp-inverse-full}
 \boxed{
 \mathscr D^{-1}(y)
 \sim
 h\!\left(
 R_0+\sum_{m=1}^\infty
 \frac{(-1)^m}{m!}
 \left(\frac{d}{dR_0}\right)^{m-1}
 \mathcal A(R_0)^m
 \right)-1.}
\end{equation}
This formula generates all pure periodic and mixed Bernoulli--periodic
coefficients.
\end{theorem}

\begin{proof}
This is the Lagrange reversion formula for the near-identity map
\(s\mapsto s+\tau\mathcal A(s)\), followed by composition with \(h\).
The parameter \(\tau\) makes the formal meaning explicit.  In the asymptotic
regime, each derivative with respect to \(R_0\) introduces factors controlled
by \(h'(R_0)=2/\Lambda\), so the resulting expansion is ordered by inverse
powers of \(\Lambda\), with additional powers of \(X^{-1}\) from \(\Omega\)
and from curvature of the core.
\end{proof}

\section{Leading oscillatory hierarchy}

Let \(X=h(R_0)\), \(\Lambda=\log X\), and abbreviate
\(P=P(X-1)\), with primes denoting derivatives of \(P\) at \(X-1\).
The leading pure-periodic displacement is
\begin{align}
 \mathscr D^{-1}(y)-(X-1)
 \sim{}&
 -\frac{2P}{\Lambda}
 +\frac{4PP'}{\Lambda^2}
 \notag\\
 &-\frac{8}{\Lambda^3}
 \left(P(P')^2+\frac12P^2P''\right)
 +\cdots,
 \label{eq:periodic-hierarchy}
\end{align}
up to mixed corrections suppressed by at least one power of \(X^{-1}\).
The first Bernoulli contribution is
\[
 -\frac{2\Omega(X)}{\Lambda}
 \sim\frac{1}{6\Lambda X}.
\]
Formula \cref{thm:phase-reversion} is preferable to hand expansion beyond
these first terms.

\chapter{Complex inverse sheets}

\section{Logarithmic and Lambert branches}

For complex \(y\), choose a logarithmic sheet
\[
 R_{\eps,\ell}:=\Log y+2\pi\ii\ell-A_\eps,
 \qquad \ell\in\mathbb Z,
\]
and a Lambert branch \(W_k\).  Define
\[
 v_{k,\ell}:=
 W_k\!\left(\frac{2R_{\eps,\ell}}{\e}\right),
 \qquad
 X_{k,\ell}:=\frac{2R_{\eps,\ell}}{v_{k,\ell}},
 \qquad
 \Lambda_{k,\ell}:=v_{k,\ell}+1.
\]
In any sector where these choices define a simple local inverse and
\(X_{k,\ell}\to\infty\), the same coefficient functions
\(d_{2n-1}(\Lambda)\) give
\[
 x_{\eps;k,\ell}(y)
 \sim
 X_{k,\ell}-1+
 \sum_{n\ge1}
 d_{2n-1}(\Lambda_{k,\ell})
 X_{k,\ell}^{-(2n-1)}.
\]
The coefficients are universal; only the sheet data change.

\section{Branch caveats}

The gamma function is meromorphic and \(\log\Gamma\) is multivalued.
Consequently, not every formal pair \((k,\ell)\) corresponds to a globally
distinct inverse branch throughout the plane.  The formula is local and
sectorial.  Branch collisions occur where \(D_\eps'(x)=0\), equivalently
\[
 \log2+\psi(x/2+1)=0.
\]
Near such points the simple inverse expansion must be replaced by a
Puiseux-type local analysis.

\chapter{Symbolic coefficient algorithms}

\section{Forward coefficients}

The forward logarithmic coefficients are direct Bernoulli expressions.
The multiplicative coefficients can be generated in \(\Ord(M^2)\) arithmetic
operations by \cref{eq:sm-recurrence}.

\begin{algorithm}[Forward coefficient table]
Given \(M\):
\begin{enumerate}[label=\arabic*.]
\item Compute \(c_k\) for \(1\le k\le\lceil M/2\rceil\).
\item Set \(\gamma_{2k-1}=c_k\), \(\gamma_{2k}=0\).
\item Set \(s_0=1\).
\item For \(m=1,\ldots,M\), compute
\[
 s_m=\frac1m\sum_{j=1}^m j\gamma_js_{m-j}.
\]
\end{enumerate}
\end{algorithm}

\section{Inverse coefficients}

A Wolfram Language implementation of the triangular recurrence is:

\begin{lstlisting}
ClearAll[c, a, z, lam, ser, e, n, k];

c[k_Integer?Positive] :=
  (1 - 2^(2 k - 1)) BernoulliB[2 k]/(2 k (2 k - 1));

a[k_Integer?Positive] := 2 c[k];

InverseBlocks[N_Integer?Positive] := Module[
  {ser = 0, coeff, en, table = <||>},
  Do[
    coeff = SeriesCoefficient[
      (1 + ser) Log[1 + ser] - ser +
      Sum[a[k] z^k (1 + ser)^(-(2 k - 1)), {k, 1, n}],
      {z, 0, n}
    ];
    en = Together[-coeff/lam];
    AssociateTo[table, n -> en];
    ser = Together[ser + en z^n],
    {n, 1, N}
  ];
  table
];
\end{lstlisting}

Here \(z=q^2\), `lam` represents \(\Lambda\), and the returned value at key
\(n\) is \(e_n=d_{2n-1}\).

\section{Coefficient-audit identities}

Any implementation should verify the following independent checks:
\begin{align}
 e_1&=-a_1/\Lambda,\label{eq:audit1}\\
 \lim_{\Lambda\to\infty}\Lambda e_n(\Lambda)&=-a_n,\label{eq:audit2}\\
 \deg_{\Lambda^{-1}}e_n&\le2n-1,\label{eq:audit3}\\
 \coeff{q^{2j}}{
 \Lambda E_N+H(E_N)+
 \sum_{k=1}^Na_kq^{2k}(1+E_N)^{-(2k-1)}
 }&=0
 \quad(1\le j\le N).
 \label{eq:audit4}
\end{align}
The last identity is a residual certificate rather than a comparison against
a precomputed table.

\chapter{Numerical validation}

\section{Forward logarithmic errors}

The following table compares the exact
\(\log D_\eps(n)\) with the centered logarithmic expansion truncated after
\(N\) Bernoulli terms.  Entries are approximation minus exact value.

\begin{table}[ht]
\centering
\small
\setlength{\tabcolsep}{3pt}
\caption{Errors in the centered logarithmic expansion.}
\label{tab:forward-errors}
\begin{tabular}{c c r r r r r}
\toprule
\(\eps\)&\(n\)&\(N=0\)&\(N=1\)&\(N=2\)&\(N=3\)&\(N=4\)\\
\midrule
0&20&
 \(3.96616{\times}10^{-3}\)&
 \(-2.09362{\times}10^{-6}\)&
 \(5.98269{\times}10^{-9}\)&
 \(-4.14413{\times}10^{-11}\)&
 \(5.30661{\times}10^{-13}\)\\
1&21&
 \(3.78606{\times}10^{-3}\)&
 \(-1.82137{\times}10^{-6}\)&
 \(4.74399{\times}10^{-9}\)&
 \(-2.99567{\times}10^{-11}\)&
 \(3.49749{\times}10^{-13}\)\\
0&100&
 \(8.25064{\times}10^{-4}\)&
 \(-1.88702{\times}10^{-8}\)&
 \(2.34020{\times}10^{-12}\)&
 \(-7.04697{\times}10^{-16}\)&
 \(3.92938{\times}10^{-19}\)\\
\bottomrule
\end{tabular}
\end{table}

The alternating signs agree with \cref{eq:RN-sign-bound}.

\section{Inverse errors}

For each test, \(y\) was generated from an exact branch value
\(y=D_\eps(x)\).  The table gives the error of the inverse expansion through
\(d_{2N-1}X^{-(2N-1)}\); \(N=0\) means the Lambert core alone.

\begin{table}[ht]
\centering
\small
\setlength{\tabcolsep}{3pt}
\caption{Errors of the parity-resolved inverse transseries.}
\label{tab:inverse-errors}
\begin{tabular}{c c r r r r r}
\toprule
\(\eps\)&\(x\)&\(N=0\)&\(N=1\)&\(N=2\)&\(N=3\)&\(N=4\)\\
\midrule
0&20&
 \(-2.60549{\times}10^{-3}\)&
 \(1.75197{\times}10^{-6}\)&
 \(-4.77514{\times}10^{-9}\)&
 \(3.10287{\times}10^{-11}\)&
 \(-3.81720{\times}10^{-13}\)\\
1&21&
 \(-2.44974{\times}10^{-3}\)&
 \(1.49560{\times}10^{-6}\)&
 \(-3.71702{\times}10^{-9}\)&
 \(2.20413{\times}10^{-11}\)&
 \(-2.47399{\times}10^{-13}\)\\
0&100&
 \(-3.57548{\times}10^{-4}\)&
 \(9.58050{\times}10^{-9}\)&
 \(-1.14691{\times}10^{-12}\)&
 \(3.31593{\times}10^{-16}\)&
 \(-1.80329{\times}10^{-19}\)\\
1&101&
 \(-3.53289{\times}10^{-4}\)&
 \(9.27850{\times}10^{-9}\)&
 \(-1.08916{\times}10^{-12}\)&
 \(3.08781{\times}10^{-16}\)&
 \(-1.64658{\times}10^{-19}\)\\
\bottomrule
\end{tabular}
\end{table}

The data validate the predicted odd inverse-power grid and show that the first
Bernoulli block removes roughly three to five decimal orders in these examples.

\chapter{Generalization to multifactorials}

\section{Residue-class branches}

The double-factorial calculation is the \(p=2\) case of a common
multifactorial normal form.  Let \(p\ge1\), and choose a residue representative
\(r\in\{1,\ldots,p\}\).  On integers \(x\equiv r\pmod p\), define
\[
 x!_{(p)}=x(x-p)(x-2p)\cdots r.
\]
Its natural residue-class continuation is
\begin{equation}\label{eq:multi-branch}
 D_{p,r}(x)
 :=
 C_{p,r}p^{x/p}\Gamma\!\left(\frac{x}{p}+1\right),
 \qquad
 C_{p,r}:=\frac{p^{1-r/p}}{\Gamma(r/p)}.
\end{equation}
All residue branches again differ only by constants.

Set
\[
 u:=x+\frac p2,
 \qquad
 A_{p,r}:=
 \log C_{p,r}+\frac12\log\!\left(\frac{2\pi}{p}\right).
\]
Then
\begin{equation}\label{eq:multi-normal}
 \log D_{p,r}(x)
 =
 \frac{u}{p}(\log u-1)+A_{p,r}+\Omega_p(u),
\end{equation}
where
\begin{align}
 \Omega_p(u)
 &\sim\sum_{k\ge1}
 \frac{c_k^{(p)}}{u^{2k-1}},\label{eq:multi-Omega}\\
 c_k^{(p)}
 &:=
 \frac{p^{2k-1}B_{2k}(\tfrac12)}
 {2k(2k-1)}.\label{eq:multi-c}
\end{align}
Its exact Borel kernel is
\begin{equation}\label{eq:multi-borel}
 \Omega_p(u)
 =
 \int_0^\infty\e^{-us}
 \frac{\frac{ps}{2\sinh(ps/2)}-1}{ps^2}\,ds.
\end{equation}
The nearest Borel poles lie at \(s=\pm2\pi\ii/p\), so the optimal forward
remainder scale is \(\e^{-2\pi u/p}\).

\section{Universal multifactorial inverse recurrence}

Let
\[
 R=\log y-A_{p,r},\qquad
 v=\W_0\!\left(\frac{pR}{\e}\right),\qquad
 X=\frac{pR}{v},\qquad
 \Lambda=v+1.
\]
Then \(X\) solves \(\frac Xp(\log X-1)=R\).  With
\(u=X(1+E)\) and \(q=X^{-1}\), the normalized equation is
\begin{equation}\label{eq:multi-master}
 \Lambda E+H(E)
 +\sum_{k\ge1}
 a_k^{(p)}q^{2k}(1+E)^{-(2k-1)}=0,
 \qquad
 a_k^{(p)}:=p\,c_k^{(p)}.
\end{equation}
Every theorem of \cref{thm:inverse-recursion,cor:Bell-recurrence,thm:lagrange-closed,thm:polynomial-structure} remains valid after replacing
\(a_k\) by \(a_k^{(p)}\).  Thus the double-factorial recurrence is one member
of a universal centered multifactorial inversion calculus.

\chapter{Conclusions}

The main structural results can be compressed into five statements.

\begin{enumerate}[label=\textup{\arabic*.},leftmargin=2.2em]
\item The shift \(u=x+1\) places both parity branches in the same centered
      gamma normal form.  Parity survives only in \(A_\eps\).
\item The universal correction has the exact Borel transform
      \[
       \widehat\Omega(s)=\frac{s/\sinh s-1}{2s^2},
      \]
      hence explicit Bernoulli/eta coefficients, a signed remainder, and a
      completely known pole lattice.
\item Exact inversion of the dominant phase gives the Lambert core
      \(X=2R_\eps/\W_0(2R_\eps/\e)\).  Corrections lie only at odd powers
      \(X^{-1},X^{-3},\ldots\).
\item Every inverse coefficient is determined by a triangular Bell recurrence
      and, independently, by a closed Lagrange--B\"urmann extraction.
      Expanding the Lambert core with Stirling cycle numbers completes the
      iterated-logarithm transseries.
\item The parity-resolved inverse, the discrete generalized inverse, and the
      inverse of a single oscillatory interpolation are different objects.
      Their relation is explicit, so no ambiguity needs to be hidden.
\end{enumerate}

The coefficient formulas are suitable for symbolic generation to arbitrary
order, while the residual bounds and exact monotonicity make them suitable for
certified numerical inversion.  The same centered mechanism extends without
essential change to every residue class of every multifactorial.

\appendix

\chapter{Coefficient tables}

\section{Logarithmic Bernoulli coefficients}

\begin{longtable}{c r r}
\caption{The coefficients \(c_k\) and \(a_k=2c_k\).}
\label{tab:c-a}\\
\toprule
\(k\)&\(c_k\)&\(a_k\)\\
\midrule
\endfirsthead
\toprule
\(k\)&\(c_k\)&\(a_k\)\\
\midrule
\endhead
1&\(-1/12\)&\(-1/6\)\\
2&\(7/360\)&\(7/180\)\\
3&\(-31/1260\)&\(-31/630\)\\
4&\(127/1680\)&\(127/840\)\\
5&\(-511/1188\)&\(-511/594\)\\
6&\(1414477/360360\)&\(1414477/180180\)\\
7&\(-8191/156\)&\(-8191/78\)\\
8&\(118518239/122400\)&\(118518239/61200\)\\
\bottomrule
\end{longtable}

\section{Two additional inverse blocks}

For reference,
{\small
\begin{align}
 d_{11}(\Lambda)
 &=-\frac{1}{5884534656000\Lambda^{11}}
 \Bigl(
 46195687238400\Lambda^{10}
 +8854370366400\Lambda^9
 +2171643318000\Lambda^8
 \notag\\
 &\quad
 +605515022112\Lambda^7
 +212869408752\Lambda^6
 +84136852800\Lambda^5
 +35589153600\Lambda^4
 \notag\\
 &\quad
 +14234019800\Lambda^3
 +4868113250\Lambda^2
 +1213962750\Lambda
 +165540375
 \Bigr),\label{eq:d11}\\[2mm]
 d_{13}(\Lambda)
 &=\frac{1}{35307207936000\Lambda^{13}}
 \Bigl(
 3707709489792000\Lambda^{12}
 +571674673451520\Lambda^{11}
 \notag\\
 &\quad
 +109005523604352\Lambda^{10}
 +22558151064384\Lambda^9
 +5893694526096\Lambda^8
 +1791313043520\Lambda^7
 \notag\\
 &\quad
 +647003557200\Lambda^6
 +251996865120\Lambda^5
 +98865967200\Lambda^4
 +35397962600\Lambda^3
 \notag\\
 &\quad
 +10534674150\Lambda^2
 +2254502250\Lambda
 +260134875
 \Bigr).
 \label{eq:d13}
\end{align}
}

\chapter{Proof of the phase reversion formula}

For completeness, introduce a formal parameter \(\tau\) and consider
\[
 Y=s+\tau A(s).
\]
Write \(s=Y+\sigma\).  Then
\[
 \sigma=-\tau A(Y+\sigma).
\]
The Lagrange--B\"urmann formula applied to
\(\sigma=\tau\phi(\sigma)\) with
\(\phi(\sigma)=-A(Y+\sigma)\) yields
\[
 \sigma
 =\sum_{m\ge1}\frac{\tau^m}{m!}
 \left.
 \frac{d^{m-1}}{d\sigma^{m-1}}
 \bigl[-A(Y+\sigma)\bigr]^m
 \right|_{\sigma=0}.
\]
Translation invariance of differentiation gives
\[
 s=Y+\sum_{m\ge1}
 \frac{(-\tau)^m}{m!}
 \left(\frac{d}{dY}\right)^{m-1}A(Y)^m.
\]
This proves \cref{eq:phase-s} and also explains why derivatives, rather than
ordinary powers alone, organize perturbations of an invertible phase.

\chapter{Notation index}

\begin{longtable}{p{0.19\textwidth}p{0.70\textwidth}}
\toprule
Symbol&Meaning\\
\midrule
\(D_\eps(x)\)&Parity-resolved continuation, \cref{eq:D-eps}.\\
\(u=x+1\)&Centered variable for the double factorial.\\
\(A_\eps\)&Parity constant: \(\tfrac12\log\pi\) or
\(\tfrac12\log2\).\\
\(\Omega(u)\)&Universal centered gamma correction,
\cref{eq:Omega-def}.\\
\(c_k\)&Forward logarithmic coefficient, \cref{eq:ck-three}.\\
\(s_m\)&Forward multiplicative coefficient,
\cref{eq:sm-complete-bell}.\\
\(R_\eps\)&Shifted target logarithm \(\log y-A_\eps\).\\
\(v\)&Lambert value \(\W_0(2R_\eps/\e)\).\\
\(X\)&Exact dominant inverse core \(2R_\eps/v\).\\
\(\Lambda\)&Core logarithm \(\log X=v+1\).\\
\(q\)&Small inverse-core parameter \(X^{-1}\).\\
\(E\)&Relative inverse correction \(u/X-1\).\\
\(e_n=d_{2n-1}\)&Coefficient of \(q^{2n}\) in \(E\), equivalently of
\(X^{-(2n-1)}\) in the inverse.\\
\(L_1,L_2\)&\(\log(2R_\eps/\e)\) and \(\log L_1\).\\
\(\widehat B_{r,m}\)&Ordinary partial Bell convolution,
\cref{eq:ordinary-Bell}.\\
\bottomrule
\end{longtable}

\backmatter

\begin{thebibliography}{99}

\bibitem{OEIS-A006882}
OEIS Foundation Inc.,
\emph{The On-Line Encyclopedia of Integer Sequences, A006882:
Double factorials \(n!!\)},
\url{https://oeis.org/A006882},
accessed September 3, 2026.

\bibitem{Meserve1948}
B. E. Meserve,
\emph{Double Factorials},
American Mathematical Monthly \textbf{55} (1948), 425--426.

\bibitem{DLMF-Gamma}
NIST Digital Library of Mathematical Functions,
\emph{Chapter 5: Gamma Function, \S5.11 Asymptotic Expansions},
\url{https://dlmf.nist.gov/5.11},
accessed September 3, 2026.

\bibitem{DLMF-Bernoulli}
NIST Digital Library of Mathematical Functions,
\emph{Chapter 24: Bernoulli and Euler Polynomials, \S24.4 Basic Properties},
\url{https://dlmf.nist.gov/24.4},
accessed September 3, 2026.

\bibitem{DLMF-W}
NIST Digital Library of Mathematical Functions,
\emph{Chapter 4: Lambert \(W\)-Function, \S4.13},
\url{https://dlmf.nist.gov/4.13},
accessed September 3, 2026.

\bibitem{Corless1996}
R. M. Corless, G. H. Gonnet, D. E. G. Hare,
D. J. Jeffrey, and D. E. Knuth,
\emph{On the Lambert \(W\) Function},
Advances in Computational Mathematics \textbf{5} (1996), 329--359,
doi:10.1007/BF02124750.

\bibitem{Comtet}
L. Comtet,
\emph{Advanced Combinatorics},
D. Reidel, Dordrecht, 1974.

\bibitem{Gessel2016}
I. M. Gessel,
\emph{Lagrange Inversion},
Journal of Combinatorial Theory, Series A \textbf{144} (2016), 212--249.

\bibitem{Nemes2015}
G. Nemes,
\emph{Error Bounds and Exponential Improvements for the Asymptotic
Expansions of the Gamma Function and Its Reciprocal},
Proceedings of the Royal Society of Edinburgh, Section A
\textbf{145} (2015), 571--596;
preprint \url{https://arxiv.org/abs/1310.0166}.

\bibitem{ParisKaminski}
R. B. Paris and D. Kaminski,
\emph{Asymptotics and Mellin--Barnes Integrals},
Cambridge University Press, Cambridge, 2001.

\bibitem{Olver}
F. W. J. Olver,
\emph{Asymptotics and Special Functions},
A K Peters, Wellesley, MA, 1997.

\bibitem{ProveIt-CCC}
V. Reshetnikov,
\emph{Combinatorial Coefficient Calculus},
ProveIt repository, semi-formalized research frontiers,
\url{https://github.com/VladimirReshetnikov/ProveIt/tree/main/Analysis/FabiusFunction/docs/semi-formalized-research-frontiers/drafts/combinatorial-coefficient-calculus/Combinatorial_Coefficient_Calculus},
accessed September 3, 2026.

\bibitem{ProveIt-Series}
V. Reshetnikov,
\emph{Series and Transseries},
ProveIt repository, semi-formalized research frontiers,
\url{https://github.com/VladimirReshetnikov/ProveIt/tree/main/Analysis/FabiusFunction/docs/semi-formalized-research-frontiers/drafts/series-and-transseries},
accessed September 3, 2026.

\bibitem{ProveIt-SpecialInverse}
V. Reshetnikov,
\emph{Special-Function Inversion at Infinity},
ProveIt repository, series-and-transseries collection,
\url{https://github.com/VladimirReshetnikov/ProveIt/tree/main/Analysis/FabiusFunction/docs/semi-formalized-research-frontiers/drafts/series-and-transseries/special-function-inversion},
accessed September 3, 2026.

\end{thebibliography}

\end{document}
